\section{Alternatives tested, limitations and conclusion}
\label{sec:discussion}

\subsection{An alternative we tested: a ``missing'' category}

The analysis left an open question (Section~\ref{sec:eda}): are the missing job
values random? On the development set, the 2{,}233 people with a missing
\code{workclass} all have a missing \code{occupation} too (2{,}241 missing
occupations in total: the other 8 people never worked), and only 9.4\% of the
people with a missing occupation earn \pos, against 23.9\% overall: they are
mostly people without a job. The most-frequent imputation hides this, since it
gives them the most common job (\code{Private}, \code{Prof-specialty}).
Imputing a separate ``missing'' category keeps the information. On the untuned
models (Table~\ref{tab:missing}, appendix), it gains 0.003 AP for the logistic
regression and 0.001 for LightGBM, with no change in F1 for LightGBM: of the
order of the variation of AP across folds. Adopting it would have meant
re-running the whole tuning for no measurable gain, so we kept the
most-frequent imputation. A separate category remains the more principled
choice for a model in production, because it does not invent a job for people
who have none.

\subsection{Overfitting: what we checked and what we did not}
\label{sec:overfitting}

\begin{itemize}
  \item \textbf{Overfitting the validation folds: checked.} The 95 Optuna
  trials and the thresholds were all chosen on the same 5 folds, so the best
  CV scores are optimistic by construction. The untouched test set measures by
  how much: every model loses about 0.01 AP (logistic regression 0.784 to
  0.773, random forest 0.814 to 0.805, LightGBM 0.832 to 0.821), and the
  ranking does not change. This is the drop expected from selection.
  \item \textbf{Memorizing the training rows: checked for the logistic
  regression only.} Its train--validation gap stays below 0.004 even with
  almost no penalty (Section~\ref{sec:regularization}). For the two tree
  models we never computed a training score. The evidence is indirect: the
  untuned random forest, whose leaves are pure, is the worst model, and
  limiting its depth and leaf size raised its CV AP by 0.064, which is what
  reducing overfitting looks like; LightGBM's 350 trees were chosen by CV
  rather than by early stopping.
\end{itemize}
For tree ensembles, a training score above the validation score is normal:
what matters is that the validation and test scores hold, and they do. A table
of train, validation and test AP for the three models, and a LightGBM learning
curve, would measure the size of the gap.

\subsection{Limitations}

\begin{itemize}
  \item \textbf{Old data.} The census dates from 1994; incomes, jobs and the
  meaning of \$50K have changed, so the model would need retraining on recent
  data.
  \item \textbf{Possible leakage.} \code{capital-gain} and \code{capital-loss}
  are declared for the same tax year as the income (Section~\ref{sec:eda}).
  \item \textbf{F1 stands in for business costs.} With real costs for each
  type of error, the threshold should minimize the expected cost instead.
  \item \textbf{Threshold chosen on fold models.} It comes from the
  out-of-fold probabilities of five models trained on 80\% of the development
  set, and is applied to a model refitted on all of it; their scores are close
  but not identical.
  \item \textbf{Probabilities not calibrated} as-is
  (Section~\ref{sec:calibration}), and overfitting of the tree models not
  measured directly (Section~\ref{sec:overfitting}).
  \item \textbf{No fairness audit.} The model uses \code{sex} and \code{race}
  as inputs, and the share of \pos{} is 30\% for men against 11\% for women.
  The audit was removed from the scope, so we did not check whether errors are
  spread equally across groups; this would be needed before any real use.
\end{itemize}

\subsection{Conclusion}

Our final model is a tuned LightGBM inside a single scikit-learn Pipeline. On
the untouched test set it reaches an AP of 0.821 and an F1 of 0.722 at its
frozen threshold of 0.54, within 0.01 of its cross-validation scores. What
mattered most was the choice of model (LightGBM is about 0.05 AP ahead of the
logistic regression), then the decision threshold, then the engineered
features. Its decisions are reliable, but its raw probabilities are
over-estimated because of the class weight, so they should be recalibrated
before being read as true chances (Section~\ref{sec:calibration}).
